\section{Numerical Implementation}
  \label{sec:numerics}

  \subsection{Computational parameters}

    We perform computations at three primes $q\in\{211,307,401\}$.
    The choice of $q=401$ extends beyond the $q=307$ plateau of O11 to confirm convergence
    of $\hat{\delta}_{\mathrm{exact}}$ in a genuine asymptotic regime.
    At each prime, the group $G_q$ has order $q^3\in\{9{,}393{,}931;\,28{,}934{,}443;\,64{,}481{,}201\}$.
    BFS is performed from the identity element; the BFS is truncated once a fraction $f$
    of all $|G_q|$ nodes has been visited.
    The precise value of $f$ at each prime is reported in Table~\ref{tab:params}.

  \subsection{Per-block span-tracking algorithm}

    For each sampled generic block $H_c$ with multi-index $c=(c_1,c_2,c_3)$, we maintain
    a $q\times q$ orthonormal basis $\{e_1,\ldots,e_k\}$ of the accumulated span
    $W_{\leq n}^{(c)}\subset\mathbb{C}^q$ using incremental Gram--Schmidt orthonormalisation.

    At each BFS step, for each new group element $g=(a,b,\gamma)\in S_n$ and each
    path $\pi=(s_1,s_2,s_3)\in S^3$, the fingerprint vector $v(\pi,g)\in G_q^3$ is
    computed, then projected onto $H_c$ via the action~\eqref{eq:block-action}.
    The projection is computed exactly: the phase $\psi_{c_i}(\gamma_i+b_i x)$ at
    each factor is evaluated modulo $q$ and accumulated without approximation.
    The resulting vector $w\in\mathbb{C}^q$ is then orthogonalised against the current
    basis; if $\|w - \mathrm{proj}_{W_{\leq n}}w\|> \varepsilon_{\mathrm{GS}}$, the
    normalised residual is added to the basis and $\Delta r_n^{(c)}$ is incremented.
    We use $\varepsilon_{\mathrm{GS}}=10^{-10}$ throughout.

    \paragraph{Computational cost.}
      The theoretical cost per BFS step is $O(q^5)$: there are $O(q^3)$ generic blocks,
      each requiring $O(q)$ operations for the Weil action and $O(q)$ for the Gram--Schmidt
      step.
      In practice, this is reduced to $O(q^4)$ per step by sampling only $M_{\mathrm{block}}$
      blocks at each prime (see Section~\ref{sec:sampling}).
      At $q=401$ with $M_{\mathrm{block}}=500$, the estimated wall-clock time per BFS step
      is of order one second on a single core.

    \paragraph{Sparsity.}
      The shell $S_n$ has size $|S_n|\sim 4n^3$ for large $n$ (derivative of the ball growth
      law $|B_n|\sim n^4$).
      For $n\leq n^*$, only a small subset of group elements is processed at each BFS step,
      so the per-step Weil computation touches a fraction $|S_n|/|G_q|\sim 4n^3/q^3\ll 1$
      of the group.

\subsection{Sampling strategy for generic blocks}
  \label{sec:sampling}

  The total number of generic multi-indices is $O(q^3)$: at $q=307$ this is approximately
  $2.9\times 10^7$, far beyond what can be computed exhaustively.
  We sample $M_{\mathrm{block}}$ blocks at each prime using the following protocol.

  \textbf{Uniform random sampling.}
  At each prime, $M_{\mathrm{block}}$ multi-indices are drawn uniformly at random from
  the set of generic multi-indices (Definition~\ref{def:generic}).
  We use $M_{\mathrm{block}}=500$ for $q\in\{211,307\}$ and $M_{\mathrm{block}}=300$
  for $q=401$ (due to the higher per-block cost).

  \textbf{Reproducibility and inter-seed dispersion.}
  For each prime, the sampling is repeated over $N_{\mathrm{seed}}=5$ independent random
  seeds.
  We report both the mean estimate $\hat{\delta}_{\mathrm{exact}}$ and the inter-seed
  standard deviation $\sigma_{\mathrm{seed}}(\hat{\delta})$, which quantifies the
  sensitivity of the extraction to the particular set of blocks sampled.
  A small $\sigma_{\mathrm{seed}}$ provides evidence that $M_{\mathrm{block}}$ blocks
  are sufficient for a representative sample.

  \textbf{Stratification.}
  As a robustness check, we additionally run a stratified sample in which $c_1,c_2,c_3$
  are drawn uniformly but with $c_1\leq c_2\leq c_3$ (up to permutation), ensuring
  coverage of different ordering classes.
  Results from the stratified and uniform samples are compared in the sensitivity
  analysis (Section~5.5).

\subsection{Fitting procedure}
  \label{sec:fitting}

  For each prime $q$ and each seed, the log-log regression
  $\log\bar{\Sigma}_n = -\delta\log n + \log C$ is performed over the fitting window
  $[n_0,n_1]$ by ordinary least squares.
  The window boundaries $n_0$ and $n_1$ are determined as follows.
  The lower bound $n_0$ is the smallest $n$ for which $|B_n|\geq q$ (ensuring that
  the BFS has covered at least one full Weil-block ambient dimension).
  The upper bound $n_1$ is the last depth before the mean saturation flag
  $\bar{\Sigma}_n < \varepsilon_{\mathrm{sat}}=10^{-3}$ is triggered.
  Sensitivity to $[n_0,n_1]$ is assessed by varying $n_0$ by $\pm 2$ and $n_1$
  by $\pm 3$ and reporting the range of resulting $\hat{\delta}$ values.

  \begin{table}[ht]
    \centering
    \caption{Computational parameters at each prime.
      $N_{\mathrm{node}}=|G_q|$; $f$ = BFS node fraction;
      $M_{\mathrm{block}}$ = blocks per seed; $N_{\mathrm{seed}}$ = number of seeds.}
    \label{tab:params}
    \begin{tabular}{ccccc}
      \hline
      $q$ & $N_{\mathrm{node}}$ & $f$  & $M_{\mathrm{block}}$ & $N_{\mathrm{seed}}$ \\
      \hline
      211 & $9.4\times 10^6$    & 0.30 & 500                  & 5                   \\
      307 & $2.9\times 10^7$    & 0.20 & 500                  & 5                   \\
      401 & $6.4\times 10^7$    & 0.10 & 300                  & 5                   \\
      \hline
    \end{tabular}
  \end{table}
